\documentclass[11pt]{article}

% Page layout
\usepackage{geometry}
\geometry{margin=1in}

% Formatting
\usepackage{xcolor}
\usepackage{setspace}
\usepackage{enumitem}
\usepackage{amssymb}
\usepackage{url}

% Boxes for reviewer comments
\usepackage{tcolorbox}
\tcbuselibrary{skins,breakable}

% Markers
\newcommand{\addressed}{\textcolor{green!60!black}{\checkmark}\ }
\newcommand{\notaddressed}{\textcolor{red!70!black}{\texttimes}\ }

\newcommand{\defin}{\mathcal D}

\newcommand{\sem}[1]{\{\![#1]\!\}}
\newcommand{\eval}[2]{#1\!\!\downarrow_{#2}}

\newcommand{\mypar}[1]{\noindent{\bf #1}}

\newcommand{\state}{\mathsf{state}}

\newcommand{\proj}{\mathsf{proj}}
\newcommand{\projE}[2]{#1\downarrow_{#2}}

\newcommand{\sconnected}{\mathsf{sConn}}
\newcommand{\hmodules}{\mathsf{hMods}}


\newcommand{\topannotation}[2]{\mathsf{top}_{#1}(#2)}

\newcommand{\marco}[1]{\textcolor{cyan}{{\footnotesize Marco: #1}}}
\newcommand{\adele}[1]{\textcolor{blue}{{\footnotesize Adele: #1}}}

\definecolor{code_indent}{HTML}{CCCCCC}

\newcommand{\indentrule}{\color{code_indent}\vrule\hspace{2pt}}

% Syntax
\definecolor{bloodyred}{rgb}{0.6,0.1,0.1}
\newcommand{\role}[1]{\textcolor{bloodyred}{\mathtt{#1}}}
%\newcommand{\role}[1]{\mathtt{#1}}
\newcommand{\nr}{\role p}
\newcommand{\nrr}{\role q}
\newcommand{\nrrr}{\role r}

\newcommand{\Var}{\mathsf{Var}}
\newcommand{\Val}{\mathsf{Val}}

\newcommand{\interactBase}[2]{{#1}\rightarrow \{{#2}\}}
\newcommand{\interactBasel}[3]{#1\rightarrow^{#2} \{#3\}}
\newcommand{\interact}[2]{\interactBase{#1}{#2}:\,\Sigma_{j\in J}\lambda_j: u_j;\ C_j}
\newcommand{\interactl}[3]{\interactBasel{#1}{#2}{#3}:\,\Sigma_{j\in J}\lambda_j: u_j;\ C_j}
\newcommand{\doubleand}{\text{\&\&}}
\newcommand{\allsynchName}{\textcolor{bloodyred}{\mathtt{allsynch}}}
\newcommand{\allsynch}[4]{\allsynchName\, \{ {\role{#1}}_{#3}:#2_{#3}\}_{#3\in #4}}

\newcommand{\command}[4] {\commandBase {#1} {#2} \Sigma_{i\in I}#3_i: #4_i}
\newcommand{\commandBase}[3] {[#1]\ #2\ \rightarrow #3}
\newcommand{\chorcommandBase}[2] {#1\ \rightarrow #2}
\newcommand{\chorcommand}[3] {\chorcommandBase {#1} \Sigma_{i\in I}#2_i: #3_i}
\newcommand{\ifTE}[4]{\textcolor{sh_keyword}{\mathtt{if}}\ #1@{#2}\ \textcolor{sh_keyword}{\mathtt{then}}\ {#3} \ \textcolor{sh_keyword}{\mathtt{else}}\ {#4}}
\newcommand{\ifTEl}[5]{\textcolor{sh_keyword}{\mathtt{if}}\ #1@^{#2}{#3} \ \textcolor{sh_keyword}{\mathtt{then}}\ {#4} \ \textcolor{sh_keyword}{\mathtt{else}}\ {#5}}
% Reply document has no listings package, so use a plain
% monospace fallback for the two code-formatting macros.
\newcommand{\codeprism}[1]{\texttt{#1}}
\newcommand{\codechor}[1]{\texttt{#1}}

% Hardcoded label definitions so that \ref{sec:...} in this reply
% reproduces the section numbers from the revised paper without
% needing the paper's .aux file.  Update these if the paper's
% numbering changes.
\makeatletter
\@namedef{r@sec:intro}{{1}{1}{Introduction}{section.1}{}}
\@namedef{r@sec:example}{{2}{1}{A Motivating Example}{section.2}{}}
\@namedef{r@sec:chor}{{3}{1}{Choreography Language}{section.3}{}}
\@namedef{r@sub:otherconstructs}{{3.3}{1}{Other language constructs}{subsection.3.3}{}}
\@namedef{r@sec:prism}{{4}{1}{The PRISM Language}{section.4}{}}
\@namedef{r@sec:proj}{{5}{1}{Projection}{section.5}{}}
\@namedef{r@def:annot}{{3}{1}{Annotation}{definition.3}{}}
\@namedef{r@thm:epp}{{1}{1}{Projection}{theorem.1}{}}
\@namedef{r@sec:bench}{{6}{1}{Benchmarking}{section.6}{}}
\@namedef{r@tab:codesize}{{1}{1}{Code size table}{table.1}{}}
\@namedef{r@ex3-res}{{3}{1}{Bitcoin PoW block-time probability}{figure.3}{}}
\@namedef{r@ass:0}{{1}{1}{Assumption}{assumption.1}{}}
\@namedef{r@sec:relwork}{{7}{1}{Related Work and Discussion}{section.7}{}}
\makeatother



\newcommand{\CEnd}{\mathbf{0}}
\newcommand{\pp}{|\!|}
\newcommand{\ppp}[1]{|[#1]|}
\newcommand{\cmd}[6]{[#1] #2 \rightarrow \{#3: #4 = #5\}_{#6}}

\newcommand{\defrec}{\stackrel{\textcolor{sh_keyword}{\mathsf{def}}}{=}}

  
% Semantics
\newcommand{\red}[1]{\ \longrightarrow_{#1}\ } 
\newcommand{\prismred}[1]{\ \stackrel{#1}{\rightsquigarrow}\ } 
\newcommand{\lab}[3]{\ [#1][#2:\;#3]\ } 

\definecolor{pblue}{rgb}{0.13,0.13,1}
\definecolor{pgreen}{rgb}{0,0.5,0}
\definecolor{pred}{rgb}{0.9,0,0}
\definecolor{pgrey}{rgb}{0.46,0.45,0.48}
\definecolor{keywordColor}{rgb}{0.50,0.00,0.33}
\definecolor{stringColor}{rgb}{0.16,0.00,1.00}
\definecolor{lineNumberColor}{rgb}{0.47,0.47,0.47}

\definecolor{carminepink}{rgb}{0.92, 0.3, 0.26}
\definecolor{jesuscolour}{rgb}{0.3, 0.92, 0.26}
\newcommand{\AV}[1]{\todo[inline, color=carminepink]{\textbf{AV}:~#1}}
\newcommand{\MC}[1]{\todo[inline, color=jesuscolour]{\textbf{MC}:~#1}}


% Very light grey background
\definecolor{reviewgray}{gray}{0.95}

\newtcolorbox{reviewercomment}{
  breakable,
  colback=reviewgray,
  colframe=black!20,
  boxrule=0.3pt,
  arc=2pt,
  left=6pt,
  right=6pt,
  top=6pt,
  bottom=6pt
}

% Commands
\newcommand{\revcomment}[1]{%
  \begin{reviewercomment}
  \textbf{Reviewer comment.} #1
  \end{reviewercomment}
}

\newcommand{\revresponse}[2][]{%
  \par\noindent\textbf{Response.} #1 #2
}

\newcommand{\revchanges}[2][]{%
  \par\noindent\textbf{Changes in the manuscript.} #1
  #2\par\bigskip
}

\begin{document}

\title{Response to Reviewers}
\author{A PROBABILISTIC CHOREOGRAPHY LANGUAGE FOR PRISM}
\date{}
\maketitle
\vspace{-0.5\baselineskip}
\begin{center}
{\large Marco Carbone and Adele Veschetti}
\end{center}

\section*{General Response}

We thank the Editor and the Reviewers for their detailed and
constructive feedback. We have revised the paper accordingly and
believe that the changes substantially improve both clarity and
technical depth.  Below we summarise the major changes and the
conventions used in this response.

\paragraph{Major structural change: projection (Section~5).}
The definition of projection has been completely rewritten. In the
previous version, the projection relied on explicit counters (natural
numbers, denoted by~$\iota$) to generate fresh identifiers.  This
mechanism was cumbersome and could be error-prone. We now adopt a more
direct approach by annotating choreographies with unique labels, which
the projection can use directly. This yields a cleaner and more easily
readable projection in the theoretical development.  For practical
reasons, our implementation still uses counters: these can be seen as
a concrete realisation of the label-generation mechanism that the
paper now abstracts into the choreography annotations. Because
Section~5 was rewritten end-to-end, its changes are \emph{not}
highlighted in blue; we kindly ask the reviewers to read it over
again.

\paragraph{Convention for this response.}
%
We have updated this response to make it more clear to the reviewers
what changes have been made with respect to the original submission.
%
Each reviewer comment is tagged with a unique identifier
(\textbf{[R$x$.$y$]}, where~$x$ is the reviewer number and~$y$ is the
sequential position of the comment). Comments are quoted literally
from the original reviews so that they can be immediately recognised.
Where the paper text was modified in response to a comment, the
corresponding change is highlighted \textbf{in blue} in the
accompanying highlighted PDF, and the same \textbf{[R$x$.$y$]} tag
appears in the margin next to the modified passage. When a single
passage in the paper addresses comments from different reviewers, the
margin tag lists both identifiers (e.g.,
\textbf{[R2.18/R3.3]}). Again, Section~5 was completely rewritten, so we did not highlight
any changes there. Minor comments (typos, spelling inconsistencies,
formatting) have been addressed throughout but are not individually
highlighted.


\newpage

% ============================================================
\section*{Reviewer 1}


\revcomment{\textbf{[R1.1]}\\
I found the technical result establishing operational equivalence
between choreographies and their projections rather weak. The result
shows that the initial step is in correspondence provided that all
projections are synchronized on some general state iota. However, it
seems to me that this synchronization is not preserved by reduction,
i.e., after a reduction step some projections remain in an iota
state while others do not. This may ultimately be irrelevant, but I
believe the issue should be discussed in the paper. I would also
expect a result establishing operational correspondence for a
sequence of reductions.
}

\revresponse{ We thank the reviewer for highlighting the mismatch
  between standard operational-correspondence results for
  choreographies and this paper focussing on PRISM.  PRISM modules are
  declarative (a fixed set of commands whose enabling depends on the
  current state), whereas choreographies have an explicit control-flow
  syntax; therefore, in the projection, control-flow is reified into
  the (extended) runtime state via the variables $s_{\nr}$, yielding
  an extended state space $S_+$. As a consequence, reductions of the
  choreography correspond to transitions on $S_+$ in the projected
  model rather than to syntactic evolution of PRISM modules.
%
  In the revised manuscript, we make this design choice explicit and
  discuss the resulting asymmetry (projection-after-reduction may
  yield fewer commands than reduction-after-projection). At the end of
  Section 5, we include a concrete example illustrating that, while
  the set of commands may differ, the projected models exhibit the
  same observable behaviour up to a suitable PRISM-induced
  equivalence; we leave the identification of the most appropriate
  equivalence (e.g., probabilistic bisimulation) and the corresponding
  strengthening of the theorem to future work.  }

\revchanges{ We expanded the discussion around Theorem 1 to clarify
  (i) why standard reduction-to-reduction correspondence does not
  directly apply to PRISM, (ii) how control-flow information is
  encoded in the extended state $S_+$ via the distributed state
  variables, and (iii) why the correspondence statement is formulated
  in terms of transitions on $S_+$ rather than syntactic changes of
  modules. We also added a detailed remark and worked example to make
  the asymmetry and its implications explicit, and we added a
  paragraph outlining how a future extension based on probabilistic
  bisimulation could capture correspondence across different command
  sets.  }

\revcomment{\textbf{[R1.2]}\\
Furthermore, the proof of this correspondence (Theorem 1) reads more
like a proof sketch than a complete proof; a more detailed and
carefully paced proof would be preferable.
}

\revresponse{ We thank the reviewer for this suggestion. In the
  revised version, we have expanded the proof of Theorem 1 and made
  its structure more explicit, in particular by spelling out the
  interaction case in detail and by making clear the key invariants on
  the extended state $S_+$ used in the correspondence argument. To
  keep the exposition readable, we do not enumerate every PRISM
  semantic step.  }

\revchanges{ Proof of Theorem 1 rewritten and expanded and interaction
  case detailed.}



\revcomment{\textbf{[R1.3]}\\
I have also some concerns about the way in which the choreography
state is projected on different participants. I think that the
presentation will benefit from a careful treatment of this aspect.
}

\revchanges{ Definitions~4 (CTMC) and~5 (DTMC) have been updated. Now
  states are pre-generated making it more clear how the projection
  works. }

% ------------------------------------------------------------

\revcomment{\textbf{[R1.4]}\\
When giving examples, the use of prime variables is not clear at
this point.
}

\revresponse{ We added a clarifying sentence in the motivating
  example paragraph of Section~1, right after the first use of
  $x'$ and~$y'$. The sentence is highlighted in blue and tagged
  \textbf{[R1.4]} in the margin.}

\revchanges{ Added the following sentence (highlighted in blue,
  tagged \textbf{[R1.4]}) in Section~1, after the description of
  the two branches:
  \begin{quote}\small
    Here, we adopt PRISM notation where the prime symbol, e.g.,
    $x'$, $y'$, indicates the next-state value of the corresponding
    variable, that is, the value it will take after the transition
    occurs.
  \end{quote}
}


% \revcomment{
% “this is the first probabilistic choreography language that is not a type abstraction”.
% What does this mean?
% }

% \revchanges{ Comment removed.  }


\revcomment{\textbf{[R1.5]}\\
I did not understand the explanation of foreach. What is the role of
k in such an expression? What is op? A formal encoding would be
welcome.
}

\revresponse{
  We apologise for the limited explanation of this construct. The basic idea is to provide an operator that quantifies over indices via the predicate $\textit{op}(i)$. The formal translation is straightforward: we add one update for each role, assuming that the relevant roles can be determined statically.
}

\revchanges{ We have updated the text where the foreach construct is
  discussed (Section~\ref{sub:otherconstructs}), with a more precise
  explanation of what it does. The new text is highlighted in blue
  and tagged \textbf{[R1.5]} in the margin.}


\revcomment{\textbf{[R1.6]}\\
I am puzzled about the notion of Non-deterministic Synchronization.
The encoding shown for the example looks quite deterministic to me.
Moreover, as written, it seems that you are lifting features of the
implementation language (PRISM) to the choreography level. The
encoding then arbitrarily introduces a synchronization. I have a
couple of questions here:
  \begin{enumerate}
  \item Does the specification determine the direction of
    communication? In the provided example, it seems that the
    interaction could also be expressed in the opposite direction,
    i.e., $q \rightarrow \{p\}$. Is the direction of the interaction
    relevant, or is it never relevant?
  \item Does any synchronization specification yield a correct
    choreography? Should there not be a requirement for some
    well-formedness condition on the allsync constructor?
  \item Could you provide a formal definition of the encoding?
  \end{enumerate}
}

\revresponse{The terminology we used is incorrect, thanks for pointing
  this out.  $\allsynchName$ is not "nondeterministic" but it is still
  a probabilistic choice. We have rewritten (and expanded) the entire
  paragraph with a clearer explanation. As for the detailed questions:
  \begin{enumerate}
    
  \item in our choreography language, writing
    $\nrr\longrightarrow\{\nr\}$ or $\nr\longrightarrow\{\nrr\}$ is
    the same. However, it has an impact on the projection, since we
    need to distinguish one of the modules (the one on the left)
    because it has a different projection.
    
  \item The notion of strong connectedness can be extended to
    $\allsynchName$, as it is, after all, an interaction between the
    modules $\nr_i$. However, this is not necessary if we first
    translate it into standard interactions (and if-then-else
    constructs), since strong connectedness can then be checked at
    that level. In the new example we have provided, strong
    connectedness does not hold because the nested if-then-else
    constructs are not connected. As this part is not a formal
    development but rather a section justifying how our implementation
    extends the theoretical framework, we do not consider it necessary
    to present a connected choreography here. In any case (as noted in
    a comment added to the projection section), these choreographies
    can be transformed into connected ones by introducing additional
    interactions.

    
  \item the translation should be straightforward once there is a
    precisely defined language for guards. See new text added. 
 \end{enumerate}
}

\revchanges{ The paragraph has been rewritten. The new text is
  highlighted in blue and tagged \textbf{[R1.6]} in the margin of
  Section~\ref{sub:otherconstructs}.}


\revcomment{\textbf{[R1.7]}\\
Example 3: The modules in this example share variables (y and x). I
wonder whether there is a choreography that describes their
behavior, since my understanding is that variables are local to
modules in choreographies.
}

\revresponse{
In the example, $x$ and $y$ are indeed owned by two different modules and are therefore shared only in read mode. In PRISM, all modules can read all variables, but only the owning module can modify its own variable. This is consistent with our choreographic reading, where updates are only performed as part of an interaction, so the write is still attributed to a specific role in that interaction. The only caveat is that, in choreographies, guards in conditionals must be local (i.e., evaluated by a single role). Thus, for this example, we could simulate the interactive execution (sync on $a$) by nesting two if-then-else constructs. However, since choreographies focus only on interactions, the internal actions (no label) cannot be simulated.
}

\revcomment{\textbf{[R1.8]}\\
The choice of annotating choreographies with labels seems somewhat
ad hoc. Why are they not introduced as a syntactic element of
choreographies (as is often done with communication channels in
standard choreographies), given that they are necessary and play a
role in the implementation?
}

\revresponse{ It is certainly true that labels are a key aspect of the
  implementation. However, the purpose of this work is to define a
  higher-level language that abstracts away from such implementation
  details. Therefore, labels are part of the projection algorithm and
  are irrelevant to the programmer who writes a
  choreography. Nevertheless, we now use an intermediate choreographic
  language and a function for annotating choreographies in a ``safe''
  way. }


\revcomment{\textbf{[R1.9]}\\
Projection of conditionals: I do not understand why the conditional
is not projected. I would expect that E should be evaluated locally
by p (hence all relevant information must be local to p). This
contrasts with the statement on page 5: ``The term if E@p then C1
else C2 denotes a system where module p evaluates the guard E
(which can contain variables located at other modules) and then
(deterministically) branches accordingly.'' I presume that there is
some underlying assumption on the projection, but I could not find
its definition.
}

\revresponse{
We believe this is a misunderstanding, and we are grateful for the opportunity to clarify it. The conditional \emph{is} projected, but only in the module of the evaluating role $\mathsf p$. This is consistent with the intended semantics: the guard $E$ is evaluated by $\mathsf p$ and the result deterministically selects the branch. The point that may be confusing is that we do \emph{not} require $E$ to be local to $\mathsf p$. As stated on page 5, $E$ may mention variables owned by other modules; in PRISM, all modules can read all variables, while only the owning module can modify them. Therefore, the projection places the conditional in $\mathsf p$'s module, and the guard can still refer to remote variables, exactly as in the PRISM syntax.
}

\revcomment{\textbf{[R1.10]}\\
Projection of choreography variable X: This relies on defs(X). I
presume that this function should be explicitly included in the
definition. Moreover, I assume that you need to impose some
conditions, such as requiring that the image of any X\_i is distinct
from nodes(C) and also distinct from the nodes in all definitions.
The passage on page 13 (``For recursive calls\ldots\ defs'') appears
somewhat rushed, and I believe defs deserves a formal definition.
}

\revresponse{
We agree with the reviewer. }

\revchanges{
In the revised formulation, counters are replaced by
explicit labels on choreographies, which makes the treatment of variables cleaner and more uniform. We have also expanded the proof for this case and
included more details to make the reasoning fully explicit.
}


\revcomment{\textbf{[R1.11]}\\
Projection of updates: Are you assuming that updates at the
choreography level are well-formed in some way? For instance, could
an assignment $y' = x$ occur where $y$ belongs to $\nr$ and $x$ to
$\nrr$? If not, this restriction should be stated clearly. Moreover,
from which element of the choreography is the assignment derived? If
I have missed this, it should nonetheless be clarified explicitly in
the paper.
}

\revresponse{
As discussed above, assignments such as $y' = x$ are permitted even when $x$ is located at $\nr$ and $y$ is located at $\nrr$. At the choreography level, the restriction is simply that, in an interaction $\interact{\nr}{\nr_1,\ldots,\nr_n}$, every update in each $u_j$ assigns to a variable owned by one of $\nr, \nr_1,\ldots,\nr_n$.
}
\revchanges{ We have added Assumption~\ref{ass:0} in
  Section~\ref{sec:chor}, highlighted in blue and tagged
  \textbf{[R1.11]} in the margin.}

\revcomment{\textbf{[R1.12]}\\
Definition 4: Is it not the case that the choreography is annotated?
Hence, the interaction $p \rightarrow \{p_1,\ldots,p_n\}$ should
also be labeled. Furthermore, from where are the labels $l_j$
drawn? Does this mean that every interaction uses the same set of
labels $l_1, \ldots$?
}

\revresponse{ Interactions are now also annotated. See new version.  }


\revcomment{\textbf{[R1.13]}\\
How are variables quantified? Are $S$, $S'$, and $\iota$ universally
quantified?
}

\revresponse{
  Yes, although now $\iota$ has disappeared.
}

\revcomment{\textbf{[R1.14]}\\
The correspondence theorem (Theorem 1) appears somewhat weak, since
it merely shows that projections can mimic the first computational
step, as far as the state is concerned. It does not provide evidence
that subsequent steps preserve equivalence. I am particularly
concerned that the projections may become misaligned with respect to
$\iota$. Thus, it is not the case that the second step necessarily
advances all projections corresponding to a given $\iota$. This
issue is clear in the case of conditionals, and I also wonder about
process variables.
}

\revresponse{ See our reply about your first comment.  }

\revcomment{\textbf{[R1.15]}\\
Although I can follow the reasoning, I wonder what the fundamental
limitation is here. Is it because the choreographies are purely
probabilistic, without non-deterministic choices? Would the dining
cryptographers require some form of non-determinism?
}

\revchanges{ The explanation is in the Dining Cryptographers paragraph
  of Section~\ref{sec:bench}, highlighted in blue and tagged
  \textbf{[R1.15]} in the margin.}

\newpage

% ============================================================
\section*{Reviewer 2}


\revcomment{\textbf{[R2.1]}\\
The main changes as stated by the authors at the end of the
Introduction are:
\begin{itemize}
\item Previously omitted details in definitions and proofs. I consider
  this a minor extension. The paper does explain which specific
  definitions, details or proofs have been included. Neither does the
  cover letter.
\item A thorough analysis of richer language constructs incorporated
  in the implementation. The Introduction does not explain which
  specific analysis and constructs they refer to. The paper does not
  specify in which section this is to be found. In principle, this
  sounds like a significant extension but with the paper does not help
  the reader in identifying it (one would need to read both papers and
  guess). Also here, the cover letter does not provide additional
  details.
\item New use cases analysed with the proposed approach. In consider
  this as a significant extension, and aligned with the recommendation
  of the COORDINATION reviews. As in the above case, the paper will be
  more clear if specific pointres to sections would be provided. As it
  is now, distinguishing the new from the old cases is left as a
  homework to the reader. My guess is that these are the
  Peer-To-Peer protocol, the leader election protocol and the dining
  cryptographers.
\end{itemize}

[Reviewer's clarification in the second round:]
``The paragraph `Changes with respect to the conference version' does
not use references to specific sections where the novel contributions
are to be found. You should add such references. Moreover, the
examples of the newly introduced additional language constructs must
be listed in the text.''
}

\revresponse{ We thank the reviewer for the clarification, and we
  apologise: in our previous reply we did not address what was
  actually being asked. We have now revised the
  ``Changes with respect to the conference version'' paragraph at the
  end of the Introduction to include explicit section references for
  each change. We kept the edit minimal: the original prose is
  preserved and only the section pointers were added inline
  (highlighted in blue); the paragraph is tagged \textbf{[R2.1]} in
  the margin.  The explicit naming of the three additional language
  constructs is addressed in the ``optional expansion'' at the end of
  this item.}

\revchanges{ In the paragraph ``Changes with respect to the
  conference version'' at the end of Section~1, we inserted the
  following explicit section pointers (each shown in blue in the
  highlighted PDF):
  \begin{itemize}
  \item ``\ldots complete proof of the main theorem
    \textbf{(Section~5)}, which were previously omitted\ldots''
  \item ``\ldots revised projection function \textbf{(Section~5)}
    that provides a cleaner formalisation\ldots''
  \item ``\ldots additional language constructs
    \textbf{(Section~3.3)}, which are also supported by the
    implementation\ldots''
  \item ``\ldots expand the related work section
    \textbf{(Section~7)}\ldots''
  \item ``\ldots new use cases analysed through choreographies
    \textbf{(Section~6)}\ldots''
  \end{itemize}

  \medskip\noindent\textbf{Optional, more verbose expansion.} The
  reviewer also asked that the three new language constructs be
  listed by name in the text. We deliberately did \emph{not} inline
  the construct names in the current revision because we felt it
  would make the paragraph cumbersome for an introduction. If the
  reviewer prefers a fully itemised and more explicit version of
  this paragraph, we are happy to replace the current one with the
  following expanded version, which also lists the new constructs
  explicitly:
  \begin{quote}\small
    This article is an expanded version of our paper presented at
    COORDINATION 2024. With respect to that version, the main
    differences are the following.
    \begin{itemize}
    \item \emph{Complete proof of the main theorem (Section~5).}
      Theorem~1, which was given only as a sketch in the conference
      version, is now accompanied by a full proof, including a
      detailed treatment of the interaction case.
    \item \emph{Revised projection function (Section~5).} The
      projection no longer relies on natural-number counters for the
      control state. Instead, choreographies are first decorated with
      explicit control and state labels through the annotation
      function $\mathsf{annot}$ (Definition~3); the projection is
      then defined directly on annotated choreographies, yielding a
      cleaner formalisation.
    \item \emph{Additional language constructs (Section~3.3).} A new
      subsection introduces three constructs that are syntactic sugar
      for the core language and are supported by our implementation:
      (i) parametric modules $\mathtt{p}[i]$, which allow indexing
      modules over a finite set; (ii) the $\mathsf{foreach}$
      construct, of the form $\mathsf{foreach}\ (k\ \mathit{op}\ i)\
      u@\mathtt{A}[k]$, which iterates an update over a set of
      indexed modules; and (iii) the $\allsynchName$ construct, which
      expresses a synchronisation in which several modules
      independently choose a branch among multiple alternatives.
    \item \emph{New case studies (Section~6).} The experimental
      evaluation includes three new case studies: the Simple
      Peer-to-Peer Protocol, the Synchronous Leader Election
      Protocol, and the Dining Cryptographers Protocol. We also added
      a table comparing the size of choreographies with that of the
      corresponding generated PRISM models.
    \item \emph{Extended Related Work and Discussion (Section~7).}
      The Related Work part has been extended with a new paragraph
      relating our setting to the implementability/realisability
      literature. The Discussion part has been expanded with a more
      precise characterisation of the class of systems amenable to
      our framework, and with a new discussion of the trade-offs of
      choreographic programming with respect to process-oriented
      programming.
    \end{itemize}
  \end{quote}
}


\revcomment{\textbf{[R2.2]}\\
For the COORDINATION paper the following remark was done in the
reviews: ``The introduction should be more clear about which kind of
concurrent process is amenable for choreographies, and for which
class of concurrent processes choreographies are not suitable.'' The
comment still applies for the new introduction.
}

\revresponse{ We added a dedicated paragraph in Section~1 that makes
  explicit which kind of concurrent systems our approach targets and
  which kinds of systems are less suited to it. The new paragraph is
  highlighted in blue and tagged \textbf{[R2.2]} in the margin of the
  revised PDF.}

\revchanges{ The following paragraph was added to Section~1, just
  before the presentation of the language (highlighted in blue and
  tagged \textbf{[R2.2]} in the margin):
  \begin{quote}\small
    Choreographic programming naturally applies to distributed
    systems whose behaviour can be described through structured
    interaction protocols. Systems whose dynamics rely heavily on
    loosely coordinated, highly unstructured, or unpredictable
    communication patterns are generally less amenable to this
    approach. Within its intended scope, however, choreographic
    programming has proved effective in practice, since
    choreographies enable the automatic generation of decentralised
    implementations that are correct by construction.
  \end{quote}
}

  
\revcomment{\textbf{[R2.3]}\\
The experimental analysis in Section 7 needs to be improved. What I
miss the most is an overview of the results, with some measurable
criteria.
}

\revresponse{ Section~\ref{sec:bench} (previously Section~7, now
  Section~6 after the removal of the old Section~6 -- see
  \textbf{[R2.4]}) has been extended with a new opening paragraph
  and a new ``Code Size Comparison'' subsection containing
  Table~\ref{tab:codesize}. Both are highlighted in blue and tagged
  \textbf{[R2.3]} in the margin. The criterion of equivalence used
  in the evaluation is clarified under \textbf{[R2.18/R3.3]}.}

\revcomment{\textbf{[R2.4]}\\
Section 6 is irrelevant and can be removed.
}

\revchanges{ Removed. The relevant information (Java/ANTLR
  implementation) was moved into the opening paragraph of
  Section~\ref{sec:bench} (tagged \textbf{[R2.3/R2.4]}). As a
  consequence, all subsequent section numbers shifted down by one.}


\revcomment{\textbf{[R2.5]}\\
It is confusing that you use lambda1,2 as label in page 2 (encoded
then as a an b), and later in page 3 lambda is the rate of the CTMC
actions. I recommend trying to make the use of symbols more coherent.
}

\revchanges{ The motivating example paragraph in Section~1 was
  rewritten so that $\lambda_1, \lambda_2$ are consistently
  presented as \emph{rates} and the labels $a, b$ are clearly
  separated as PRISM command labels. The new text is highlighted in
  blue and tagged \textbf{[R2.5]} in the margin.}

\revcomment{\textbf{[R2.6]}\\
Relate to this: the text uses somtimes the math symbols and
sometimes the full name of the greek laters lambda, mu, theta. If
you refer to code snippets, where the latter is used, I would use
that in the main text too to ease readability. Note that you refer
to other code symbols using the code format (e.g.\ CheckOutState).
}

\revchanges{ In Section~\ref{sec:example}, $\lambda, \mu, \theta$
  replaced by code-formatted \codeprism{lambda}, \codeprism{mu},
  \codeprism{theta} (two locations, highlighted in blue, tagged
  \textbf{[R2.6]}).}

\revcomment{\textbf{[R2.7]}\\
I don't know if this intentional, but if one wants to specify no
update, one has to write $x' = x$ (a void update). In PRISM one can
simply write true, if I recall properly. This is related also to
the syntax of PRISM introduced in 4.1.
}

\revresponse{
  From a theoretical point of view, there is no difference, therefore we will leave this as it is. Our implementation allows also True. 
}

\revcomment{\textbf{[R2.8]}\\
Are you sure Def.~1 works as intended? What happens if $x$ appears
in $u$? For example, $[x' = 0, y' = x]$. In Prism, $y$ will be
updated to the value of $x$ in $S$ (see remark at the bottom of
\url{https://www.prismmodelchecker.org/manual/ThePRISMLanguage/Commands}),
but in your case it will be set to the new value of $x$ (0). This
is also relevant to the definition of $\mu$ in page 9 and needs to
be fixed.
}

\revresponse{ Yes, our definition was incorrect. }

\revchanges{ Definition~1 in Section~\ref{sec:chor} has been
  rewritten to handle multiple updates simultaneously. The new
  definition is highlighted in blue and tagged \textbf{[R2.8]} in
  the margin.}

\revcomment{\textbf{[R2.12]}\\
General remark: the way you define unique identifiers for the
control points in the choreographies is not self-explainable.
Readers familiar with building control flow graphs and similar
graphical structures from code will get it, but other readers may
find it cryptic. I would suggest enriching the explaination of how
control points sq are correctly created.
}

\revchanges{ Section~\ref{sec:proj} has been rewritten from scratch
  with a new label-based mechanism (Definitions~2
  and~\ref{def:annot}) that replaces the old counter-based
  construction and makes the generation of control-point
  identifiers explicit. We kindly ask the reviewer to read
  Section~\ref{sec:proj} afresh.}


\revcomment{\textbf{[R2.13]}\\
``The projected update $u_j \downarrow q$ acts as a filter on the
list of updates in $u_j$'' -- is this defined in the paper?
}

\revchanges{ Section~\ref{sec:proj} has been rewritten from scratch.
  The explanation of $\projE{u_j}{\role q}$ now appears immediately
  after the first projection rule, together with
  Assumption~\ref{ass:0} in Section~\ref{sec:chor} (see
  \textbf{[R1.11]}). We kindly ask the reviewer to read
  Section~\ref{sec:proj} afresh.}


\revcomment{\textbf{[R2.18]}\\
``we demonstrate that PRISM behaves similarly'' -- be more precise
about what you mean by ``similarly''. Some notion of bisimilation for
CTMCs/DTMCs? Equi-satisfaction of some subset of properties? You may
want to try out STORM's bisimulation support. It could give you
higher confidence than just checking that the models behave
similarly for some specific properties.
}

\revresponse{ We apologise for not having addressed this comment in
  our previous reply. We rephrased ``behaves similarly'' as
  ``exhibits equivalent behaviour'' and clarified that we mean
  equi-satisfaction on the benchmark properties (reachability,
  time-bounded probabilities, expected values), not bisimilarity.
  We did not use STORM's bisimulation support because our
  comparison is property-based.}

\revchanges{ In Section~\ref{sec:bench}, the introductory paragraph
  was rewritten: item~(ii) now reads ``PRISM exhibits equivalent
  behaviour\ldots'' followed by a new sentence defining the
  criterion. Highlighted in blue and tagged \textbf{[R2.18/R3.3]}
  in the margin.}


\revcomment{\textbf{[R2.19]}\\
``Any discrepancies between the original and generated models are
due to inherent variations in the simulation-based calculation of
probability.'' -- have you tried not to use simulations but to use
exact calculations? Perhaps also with rationals instead of floating
points.
}

\revresponse{ We apologise for not having addressed this comment
  in our previous reply. For the four smaller benchmarks, PRISM
  performs exact model checking; simulation is used only for
  Bitcoin~PoW and Hybrid~Casper, where the state space is too large
  for exact verification. The quoted sentence refers specifically
  to those two cases. We have now made this distinction explicit in
  the paper.}

\revchanges{ Added a clarifying sentence in the preamble of
  Section~\ref{sec:bench} stating which benchmarks use exact model
  checking and which use simulation. Highlighted in blue and tagged
  \textbf{[R2.19]} in the margin.}


% (Continue similarly for remaining detailed comments as needed.)

\newpage

% ============================================================
\section*{Reviewer 3}


\revcomment{\textbf{[R3.1]}\\
The pros and cons of choreographic programming as opposed to
process-oriented programming are not elaborated upon very much. I
understand that flexibility is sacrificed in favour of
predictability with the advantage of guaranteeing
correctness-by-construction. What is lost by doing so? And while
the authors mention in the Introduction that their ``approach can
be used for a proper subset of distributed systems'', the size of
this subset and type of systems populating it is not even hinted
at. The authors only write that ``Some of the case studies
presented in the PRISM documentation [2] cannot be modeled'' and
then provide two concrete examples of what cannot be modelled.
Could the authors state something more general?
}

\revchanges{ We added a new paragraph in
  Section~\ref{sec:relwork} (Discussion) characterising what is
  lost by adopting the choreographic discipline (highlighted in
  blue, tagged \textbf{[R3.1]}). A complementary paragraph in
  Section~1 characterises which systems our approach targets (tagged
  \textbf{[R2.2]}).}




\revcomment{\textbf{[R3.2]}\\
How does this relate to the notions of implementability and
realisability as studied for transition systems and (team) automata?
In the following papers from the literature, a global transition
system or (team) automaton is given and the main question is whether
or not there exists a distributed set of local (cooperating,
synchronously communicating) transition systems or (component)
automata, whose composition is equivalent or bisimilar to the global
one.
\begin{itemize}
\item I.~Castellani, M.~Mukund, and P.S.~Thiagarajan, Synthesizing
  Distributed Transition Systems from Global Specifications. In:
  FSTTCS 1999. LNCS 1738, pp.~219--231. Springer, 1999.
\item M.~Mukund, From Global Specifications to Distributed
  Implementations. In: Synthesis and Control of Discrete Event
  Systems. Springer, 2002.
\item M.H.~ter Beek, R.~Hennicker, and J.~Proen\c{c}a, Realisability
  of Global Models of Interaction. In: ICTAC 2023. LNCS 14446,
  Springer, 2023, pp.~236--255.
\item M.H.~ter Beek, R.~Hennicker, and J.~Proen\c{c}a, Team Automata:
  Overview and Roadmap. In: COORDINATION 2024. LNCS 14676, Springer,
  2024, pp.~161--198.
\end{itemize}
}

\revchanges{ Added a new paragraph in Section~\ref{sec:relwork}
  (Related Work) relating our approach to the
  implementability/realisability literature and citing the four
  suggested references. Highlighted in blue and tagged
  \textbf{[R3.2]} in the margin.}

\revcomment{\textbf{[R3.3]}\\
Finally, during the experimental evaluation, the models generated
by the authors' approach are only compared to the original PRISM
benchmarks with respect to certain properties and functionality,
but not semantically. Could the authors imagine a more thorough
comparison, e.g., whether or not the models are bisimilar, and why?
}

\revresponse{ We apologise that our previous reply pointed to
  ``the end of Section~5'', which was misleading. The relevant
  discussion is in the introductory paragraph of
  Section~\ref{sec:bench}, where we clarify that we compare models
  on quantitative properties, not bisimilarity. We do not aim for
  semantic equivalence at the model level.}

\revchanges{ See the paragraph tagged \textbf{[R2.18/R3.3]} in
  Section~\ref{sec:bench}.}

\end{document}
